%% ma: L11-B21-0009
%% lop: 11
%% chuong: 6
%% bai: 21
%% dang: 11.21.7
%% loai: DS
%% muc: [NB, NB, TH, VD]
%% dap_an: "ĐSĐĐ"
%% nguon: "(NBV)-ĐỀ SỐ 9-MỖI NGÀY 1 ĐỀ THI 2025.pdf (D09, MỖI NGÀY 1 ĐỀ - Đề số 9), Phần II Câu 2"
%% kien_thuc: "Bài 18 (lũy thừa với số mũ thực), Bài 20 (hàm số mũ), Bài 21 (bất phương trình mũ và lôgarit)"
%% trang_thai: cho-duyet
%% ly_do: "Dạng mới đề xuất 11.21.7 (cường độ ánh sáng giảm theo độ sâu I = I0.a^d: hằng số a, giá trị, bất phương trình mũ) chờ giáo viên duyệt"
%% ghi_chu: "Đáp án gốc Đ S Đ Đ đúng (kiểm bằng sympy: a^20 = 0,3585; log_{0,95}0,5 = 13,51). Đề gốc có ảnh minh họa ánh sáng dưới biển (ảnh trang trí, không chứa dữ kiện toán) nên bỏ hình."
\begin{ex}
Cường độ ánh sáng $I$ dưới mực nước biển giảm dần theo độ sâu được biểu diễn bởi công thức $I=I_0\cdot a^d$, trong đó $I_0$ là cường độ ánh sáng tại mặt biển, $a$ là hằng số ($a>0$) và $d$ là độ sâu tính bằng mét tính từ mặt biển. Biết rằng cường độ ánh sáng tại độ sâu $1$ m bằng $0{,}95I_0$.
\choiceTF
{\True Ta có $0<a<1$.}
{Giá trị của $a=1{,}2$.}
{\True Tại độ sâu $20$ m, cường độ ánh sáng bằng $36$ phần trăm so với $I_0$ (làm tròn đến chữ số thập phân thứ hai).}
{\True Cường độ ánh sáng giảm đi quá một nửa ở độ sâu trên $14$ m dưới mực nước biển.}
\loigiai{Với $d=1$: $I_0a=0{,}95I_0$ nên $a=0{,}95$.

a) Đúng ($0<0{,}95<1$). b) Sai ($a=0{,}95$).

c) $I(20)=I_0\cdot0{,}95^{20}\approx0{,}3585I_0\approx0{,}36I_0$, tức $36\%$. Đúng.

d) $I<\dfrac12I_0\Leftrightarrow0{,}95^d<0{,}5\Leftrightarrow d>\log_{0{,}95}0{,}5\approx13{,}51$ (cơ số nhỏ hơn $1$). Vậy tại độ sâu trên $14$ m ($>13{,}51$ m) cường độ ánh sáng đã giảm quá một nửa. Đúng.}
\end{ex}
